% Kapitel 1 -- Von dem Geschlecht der Ottensteiner und Kuonrâts Herkunft

\section{Von dem Geschlecht der Ottensteiner}

\mylettrine{H}{och} über dem Kamptal, eingebettet in den \gls{Waldviertel}, erhebt sich die Burg Ottenstein auf einem Felsrücken über dem Tal. Wer sie zum ersten Male erblickt, mag sie für unbedeutend halten; wer sie jedoch kennt und ihre Geschichte versteht, der weiß, dass hier ein Geschlecht über Generationen seine Wurzeln schlug, das weder durch besonderen Glanz noch durch besonderen Ruhm, wohl aber durch unerschütterliche Beständigkeit sich hervortat.

Ältester namentlich greifbarer Burgherr war \gls{Hugo von Ottenstaine}, wie mir aus Urkunden des Stiftes \gls{Zwettl} bekannt ist: er ist bezeugt in den Jahren des Herrn 1177 und 1192, ein wohlhabender Mann, der dem Stift Schenkungen machte und dessen Burg als freies Eigen der Familie galt, nicht als herzogliches \gls{lehen}. Die Herren von Ottenstein, so heißt es, waren verwandt mit den Herren von \gls{Rauheneck} bei \gls{Baden}, und ihr Einfluss reichte tief in die Waldlandschaft des nördlichen \gls{Oesterriche}s.

Aus der Nachkommenschaft \gls{Hugo von Ottenstaine}s gingen mehrere Söhne hervor, deren Namen mir nicht alle überliefert sind. Die Hauptlinie verblieb auf der Burg Ottenstein und führte das volle Prädiktat weiter; eine jüngere Linie jedoch, aus der \gls{Hadmar von Steinare} hervorging, nannte sich schlicht von Steinære. So war es unter dem Adel im \gls{Waldviertel} üblich: der Erstgeborene erbt den Stammsitz, die Jüngeren richten sich anderwo ein, verkürzen den Namen und begnügen sich mit dem, was bleibt. \gls{Hadmar von Steinare} verwaltete seine Güter, zahlte den Zehnt und lebte, soweit ich ermitteln konnte, ein stilles, geordnetes Leben im Schatten der großen Burg, die er nie sein Eigen nennen durfte.
